\begin{abstract}
An agent cannot adapt well to a user merely by answering the current request; it must also decide what it still needs to learn about that user. This creates a measurement problem. Static Theory-of-Mind benchmarks test inference from supplied evidence, whereas interactive benchmarks typically score downstream task or social outcomes. The former removes evidence acquisition from the tested policy, and the latter can reward success without an accurate user model.

We argue that functional user modeling should instead be evaluated as active latent-state identification: the target state must be explicit, acquiring evidence must consume a limited budget, and success must be attributable to visible interaction rather than hidden-state leakage or a favorable prior. \knowact instantiates these requirements for knowledge-grounded interaction. A public domain graph defines a shared hypothesis space; a hidden user knowledge map defines the reconstruction target; and a tested agent chooses diagnostic questions before submitting a full-map estimate.

The benchmark separates simulator, agent, and evaluation contexts to address the tension between natural interaction and experimental control. A de-identified answer blueprint constrains what the simulator may express, while deterministic node-level scoring avoids an evaluator-model confound. The initial domain is statistical learning, grounded in an authoritative textbook. This paper develops the benchmark's measurement argument and validation protocol; empirical validation remains future work.
\end{abstract}
